\uuid{oiwk}
\exo7id{6066}
\titre{exo7 6066}
\auteur{queffelec}
\organisation{exo7}
\datecreate{2011-10-16}
\isIndication{false}
\isCorrection{false}
\chapitre{Continuité, uniforme continuité}
\sousChapitre{Continuité, uniforme continuité}
\cpge{mpsi-mp2i.fonctions-valeurs-reelles-complexes.a, pcsi.techniques-calcul-fonctions.a, ptsi.techniques-calcul-fonctions.a, mpsi-mp2i.limites-continuite.a, pcsi.limites-continuite.a, ptsi.limites-continuite.a, bcpst1.fonctions-reelles-usuelles.a, bcpst1.limites-continuite.c}
\module{Topologie}
\niveau{L3}
\difficulte{}

\contenu{
\question{
Soit $f$ une isométrie de $\R$ dans $\R$. Montrer qu'on a soit
 $f(x)=a-x$, soit $f(x)=a+x$, où $a=f(0)$. (Se ramener à $a=0$.)
}
}
